% Options for packages loaded elsewhere
\PassOptionsToPackage{unicode}{hyperref}
\PassOptionsToPackage{hyphens}{url}
\PassOptionsToPackage{dvipsnames,svgnames,x11names}{xcolor}
\documentclass[
  danish,
]{article}
\usepackage{xcolor}
\usepackage[margin=2.5cm]{geometry}
\usepackage{amsmath,amssymb}
\setcounter{secnumdepth}{-\maxdimen} % remove section numbering
\usepackage{iftex}
\ifPDFTeX
  \usepackage[T1]{fontenc}
  \usepackage[utf8]{inputenc}
  \usepackage{textcomp} % provide euro and other symbols
\else % if luatex or xetex
  \usepackage{unicode-math} % this also loads fontspec
  \defaultfontfeatures{Scale=MatchLowercase}
  \defaultfontfeatures[\rmfamily]{Ligatures=TeX,Scale=1}
\fi
\usepackage{lmodern}
\ifPDFTeX\else
  % xetex/luatex font selection
\fi
% Use upquote if available, for straight quotes in verbatim environments
\IfFileExists{upquote.sty}{\usepackage{upquote}}{}
\IfFileExists{microtype.sty}{% use microtype if available
  \usepackage[]{microtype}
  \UseMicrotypeSet[protrusion]{basicmath} % disable protrusion for tt fonts
}{}
\makeatletter
\@ifundefined{KOMAClassName}{% if non-KOMA class
  \IfFileExists{parskip.sty}{%
    \usepackage{parskip}
  }{% else
    \setlength{\parindent}{0pt}
    \setlength{\parskip}{6pt plus 2pt minus 1pt}}
}{% if KOMA class
  \KOMAoptions{parskip=half}}
\makeatother
\ifLuaTeX
\usepackage[bidi=basic,shorthands=off]{babel}
\else
\usepackage[bidi=default,shorthands=off]{babel}
\fi
\ifLuaTeX
  \usepackage{selnolig} % disable illegal ligatures
\fi
\setlength{\emergencystretch}{3em} % prevent overfull lines
\providecommand{\tightlist}{%
  \setlength{\itemsep}{0pt}\setlength{\parskip}{0pt}}
\usepackage{bookmark}
\IfFileExists{xurl.sty}{\usepackage{xurl}}{} % add URL line breaks if available
\urlstyle{same}
\hypersetup{
  pdflang={da-DK},
  colorlinks=true,
  linkcolor={Maroon},
  filecolor={Maroon},
  citecolor={Blue},
  urlcolor={Blue},
  pdfcreator={LaTeX via pandoc}}

\author{}
\date{}

\begin{document}

\section{Hændelsesdetektering}\label{huxe6ndelsesdetektering}

Når der registreres en lydhændelse i mikrofonerne, skal der afgøres om
lyden kom fra en drone eller ej, samt hvor den kom fra. Hvor lyden kom
fra bestemmes ved hjælp af TDOA (henvis til TDOA analyse). Det vil sige
at hændelsesdetekteringen skal afgøre om der kom en lyd fra en
``drone'', samt hvornår lyden ankom til hver mikrofon. Dette er altså de
to problemer som skal løses:

\begin{enumerate}
\def\labelenumi{\arabic{enumi}.}
\item
  At afgøre om lyden er kommet fra en drone
\item
  At bestemme hvornår lyden ankom til hver mikrofon
\end{enumerate}

\subsection{At bestemme om lyden er kommet fra en
drone}\label{at-bestemme-om-lyden-er-kommet-fra-en-drone}

Lyden kommer løbende fra vores datastream (henvis til datastream) i en
ringbuffer der konstant opdateres. Til hændelsesdetekteringen tilgås
bufferen via dens read pointer. Hvert index i ringbufferen indeholder en
frame, som består af ét sample fra hver af mikrofonnodens to mikrofoner,
og hver frame har et tilhørende timestamp. Vi sampler med 48 kHz og
bruger ALSA-formatet S32\_LE, dvs. hvert sample er et 32-bit signed
heltal (int32\_t).

Ifølge databladet for SPH0645LM4H-B leverer mikrofonen data i et 24-bit
format, men kun med 18 bits præcision, resten af bitsene er nuller. I en
32-bit container ligger de 18 gyldige bits øverst, og de nederste 14
bits er altid 0. Vi skifter derfor hvert sample 14 bits til højre, før
vi bruger det:

\[x = x_{\text{rå}} \gg 14\]

Et aritmetisk skift bevarer fortegnet, så hvert sample får en talværdi
mellem \(-131072\) og \(131071\). At de nederste 14 bits altid er 0, skal
verificeres på den faktiske opstilling.

Mikrofonen kan desuden have en DC-offset, altså en konstant forskydning
af signalet. En DC-offset tæller med i RMS-værdien og skal derfor
fjernes, inden vi kvadrerer. Det gør vi ved løbende at trække
gennemsnittet af signalet fra hvert sample. Om det er nødvendigt,
tjekkes ved at optage stilhed og se om middelværdien er tæt på 0.

Hændelsesdetekteringen står ikke selv for timestamps. Vi antager, at de
timestamps vi modtager fra ringbufferen er i en tidsbase, der er
synkroniseret mellem de to mikrofonnoder via PTP, så timestamps fra de
to noder kan sammenlignes direkte i master RPI'en.

\subsubsection{PCM tærskel}\label{pcm-tuxe6rskel}

Vi vil registrere lyden som et klap, hvis den overstiger en bestemt
tærskel i dB. Tærsklen er endnu ikke endeligt fastsat. Kravet F1 nævner
75 dB, men da lydniveauet af et klap er lavere ved de mikrofoner der
ligger længere væk, kommer tærsklen til at blive sænket, så et klap kan
detekteres alle steder inde i 3x3 meter målområdet. Det er dyrt og
unødvendigt at omregne for hver rå værdi, så derfor skal der på forhånd
udregnes en værdi i PCM-data for vores tærskel.

Vi starter med at finde fuld skala for den digitale repræsentation:

\[A_{\mathrm{fullscale}} = 2^{(N - 1)} - 1\]

Her er N antallet af bits i mikrofonens output. Det afhænger både af
mikrofonen og hvordan driveren pakker vores samples. Da vi skifter vores
samples ned til 18 bit, er N = 18.

Herefter vil vi finde den rå amplitude ved databladets reference
lydtryk:

\[A_{\mathrm{ref}} = A_{\mathrm{fullscale}} \cdot 10^{\frac{S_{\mathrm{dBFS}}}{20}}\]

Her er \(S_{\mathrm{dBFS}}\) mikrofonens sensitivity spec fra databladet. For
SPH0645LM4H-B er den typisk -26 dBFS, men den kan variere mellem -29 og
-23 dBFS fra mikrofon til mikrofon.

Nu vil vi finde den rå PCM-tærskel svarende til vores tærskel i dB:

\[\mathit{THRESHOLD}_{\mathrm{RAW}} = A_{\mathrm{ref}} \cdot 10^{\frac{\mathit{SPL}_{\mathrm{target}} - \mathit{SPL}_{\mathrm{ref}}}{20}}\]

Her er \(\mathit{SPL}_{\mathrm{target}}\) den tærskel i dB vi bestemmer, og \(\mathit{SPL}_{\mathrm{ref}}\)
er det lydtryksniveau sensitiviteten er målt ved, som for SPH0645LM4H-B
er 94 dB SPL.

Disse 3 skridt kan samles i en formel:

\[\mathit{THRESHOLD}_{\mathrm{RAW}} = (2^{N - 1} - 1) \cdot 10^{\frac{S_{\mathrm{dBFS}} + (\mathit{SPL}_{\mathrm{target}} - \mathit{SPL}_{\mathrm{ref}})}{20}}\]

Da sensitiviteten kan variere ±3 dB, kan to mikrofoner give op til 6 dB
forskel for det samme lydtryk. Det skal vi tage højde for, når tærsklen
fastsættes, f.eks. ved at måle på hver mikrofon.

\subsubsection{Tjek om lyd er oversteget
tærskel}\label{tjek-om-lyd-er-oversteget-tuxe6rskel}

Da et klap er en meget kort transient, skal RMS-vinduet ikke være alt
for stort, vi vil gerne have et 2-5 ms vindue til at vurdere om en
hændelse har fundet sted:

Antallet af samples i vinduet findes ud fra vinduets længde og vores
sample rate \(f_{s}\) = 48000 Hz:

\[N_{\mathrm{vindue}} = t_{\mathrm{vindue}} \cdot f_{s}\]

\[N \approx 2\text{--}5\ \mathrm{ms} \Longrightarrow \text{ca. } 96\text{--}240 \text{ samples ved } 48\ \mathrm{kHz}\]

Vi beregner RMS-værdien for hver af nodens to mikrofoner hver for sig. I
stedet for at dele bufferen op i faste blokke af f.eks. 200 samples,
bruger vi et glidende vindue. Hvis vi brugte faste blokke, kunne et klap
blive delt over to blokke, så ingen af dem overstiger tærsklen. Med et
glidende vindue flytter vinduet sig ét sample ad gangen sammen med read
pointeren:

\[\lbrack \mathit{readpointer} - N + 1,\ \mathit{readpointer}\rbrack\]

For hvert nyt sample kvadrerer vi værdien og lægger den til summen, mens
vi trækker kvadratet af det ældste sample i vinduet fra. På den måde
behøver vi ikke at summere hele vinduet igen for hvert sample:

\[\mathit{sum}_{\mathrm{ny}} = \mathit{sum}_{\mathrm{gammel}} + x_{n}^{2} - x_{n - N}^{2}\]

Da den kvadrerede værdi af hvert sample kan blive ekstremt høj:

Den største værdi et 18-bit sample kan have er \(2^{N - 1} - 1 = 131071\), så det største kvadrat er:

\[x_{\max}^{2} = \left( 2^{N - 1} - 1 \right)^{2}\]

\[131071^{2} = 17179607041 \approx 1{,}717961 \cdot 10^{10}\]

Skal summens type kunne indeholde et endnu højere tal:

\[\mathit{sum}_{\max} = N_{\mathrm{vindue}} \cdot \left( 2^{N - 1} - 1 \right)^{2}\]

\[17179607041 \cdot 240 = 4123105689840 \approx 4{,}123106 \cdot 10^{12}\]

Vi vælger derfor at repræsentere summen med en int64\_t, som kan
indeholde værdier op til ca. \(9{,}2 \cdot 10^{18}\). Det er også en af grundene til,
at vi skifter vores samples ned til 18 bit. Kvadratet af et råt 32-bit
sample kan blive op til ca. \(4{,}6 \cdot 10^{18}\), og allerede to af disse vil
overskride grænsen for en int64\_t.

Til sidst sammenligner vi med vores tærskel. I stedet for at udregne
RMS-værdien med en kvadratrod, sammenligner vi summen direkte med den
kvadrerede tærskel ganget med antallet af samples i vinduet. Det giver
samme resultat, men sparer en kvadratrod for hvert sample:

\[\frac{\sqrt{\mathit{sum}}}{\sqrt{N}} > \mathit{THRESHOLD}_{\mathrm{RAW}} \quad \Longleftrightarrow \quad \mathit{sum} > N \cdot \mathit{THRESHOLD}_{\mathrm{RAW}}^{2}\]

\subsubsection{Hvornår har en mikrofonnode detekteret en
hændelse}\label{hvornuxe5r-har-en-mikrofonnode-detekteret-en-huxe6ndelse}

En mikrofonnode registrerer kun en hændelse, hvis begge dens mikrofoner
overstiger tærsklen. De to mikrofoner sidder i nabohjørner af
målområdet, altså 3 meter fra hinanden. Den største forsinkelse mellem
de to mikrofoner er:

\[\Delta t_{\max} = \frac{d}{c}\]

hvor d er afstanden mellem de to mikrofoner og c er lydens hastighed,
ca. 343 m/s ved 20 °C:

\[\frac{3\ \mathrm{m}}{343\ \mathrm{m/s}} \approx 8{,}7\ \mathrm{ms}\]

Lyden kan altså ankomme op til 8,7 ms senere til den ene mikrofon end
til den anden. Den anden mikrofon skal derfor overstige tærsklen senest
ca. 8,7 ms plus en lille margin efter den første, for at det tælles som
samme hændelse. Tidspunktet hvor den første mikrofon overstiger
tærsklen, kalder vi triggerpunktet.

Efter en hændelse ignorerer detektoren nye hændelser i en periode
(holdoff). Ellers kan efterklang og ekko i rummet give flere hændelser
fra det samme klap. Holdoff skal mindst være lige så lang som den del af
udsnittet, der ligger efter triggerpunktet (se næste afsnit).

\subsection{At bestemme hvornår lyden ankom til hver
mikrofon}\label{at-bestemme-hvornuxe5r-lyden-ankom-til-hver-mikrofon}

RMS-vinduet fortæller kun, at der er sket en hændelse inden for de
seneste 2-5 ms, ikke præcis hvornår lyden ankom. Det præcise
ankomsttidspunkt bestemmes først i master RPI'en ved krydskorrelation
(henvis til TDOA analyse). Mikrofonnodens opgave er derfor at sende et
udsnit af de rå samples omkring triggerpunktet videre til master, for
begge mikrofoner.

\subsubsection{Udsnittets størrelse}\label{udsnittets-stuxf8rrelse}

Udsnittet skal indeholde starten af klappet på begge mikrofoner, samt
nok af klappet til at krydskorrelationen kan bruge det. Vi vælger et
udsnit på 50 ms, fordelt således:

\begin{itemize}
\item
  15 ms før triggerpunktet. Klappet starter før RMS-værdien overstiger
  tærsklen, da RMS-vinduet først skal fyldes, og det tager længere tid
  ved en mikrofon der ligger langt væk.
\item
  35 ms efter triggerpunktet. Det dækker de op til 8,7 ms før lyden når
  nodens anden mikrofon, de første ca. 10 ms af klappet som
  krydskorrelationen skal bruge, samt en margin.
\end{itemize}

\[N_{\mathrm{frames}} = t_{\mathrm{udsnit}} \cdot f_{s}\]

\[0{,}050\ \mathrm{s} \cdot 48000\ \mathrm{Hz} = 2400 \text{ frames}\]

Med 2 kanaler og 4 bytes pr. sample bliver det \(2400 \cdot 2 \cdot 4 = 19200\) bytes,
altså ca. 19 KB pr. hændelse. Fordelingen på 15 ms før og 35 ms efter er
et udgangspunkt og bør efterprøves ved at klappe i det fjerneste hjørne
af målområdet og se på dataene. Derfor laves begge værdier som
parametre, der kan justeres uden at ændre koden.

Når klappet detekteres, er data efter triggerpunktet endnu ikke skrevet
til ringbufferen. Detektoren venter derfor, til write pointeren er
mindst 35 ms foran triggerpunktet, før udsnittet kopieres ud af
ringbufferen. Efter kopien tjekkes det, at starten af udsnittet ikke er
blevet overskrevet i mellemtiden. Er den det, kasseres hændelsen.

\subsubsection{Afsendelse til master}\label{afsendelse-til-master}

Udsnittet sendes til master RPI'en over en TCP-forbindelse, med én fast
forbindelse fra hver mikrofonnode. TCP sikrer at hele udsnittet når frem
i den rigtige rækkefølge, hvilket UDP ikke gør, og forsinkelsen er
ubetydelig i forhold til kravet om at positionen vises senest 1 s efter
hændelsen (P3). Hver besked består af en header efterfulgt af de rå
samples:

\begin{itemize}
\item
  Node-id: hvilken mikrofonnode beskeden kommer fra
\item
  Løbenummer for hændelsen
\item
  Sample rate (48000 Hz)
\item
  Antal kanaler (2)
\item
  Antal frames i udsnittet (2400)
\item
  Timestamp for udsnittets første frame
\item
  Placeringen af triggerpunktet i udsnittet
\item
  De rå samples, hvor de to kanaler ligger skiftevis (interleaved)
\end{itemize}

Mikrofonnoden skal automatisk genoprette forbindelsen, hvis den mistes
(R2).

\subsubsection{Behandling i master}\label{behandling-i-master}

Master RPI'en parrer hændelser fra de to mikrofonnoder ud fra deres
timestamps. To hændelser hører sammen, hvis deres triggerpunkter ligger
tættere på hinanden end den største mulige tidsforskel i målområdet,
plus usikkerheden på synkroniseringen:

\[\Delta t_{\max} = \frac{d}{c}, \qquad d = \sqrt{3^{2} + 3^{2}} = 3\sqrt{2}\ \mathrm{m}\]

\[\frac{3\sqrt{2}\ \mathrm{m}}{343\ \mathrm{m/s}} \approx 12{,}4\ \mathrm{ms}\]

Hvis kun én mikrofonnode har detekteret en hændelse, kasseres den.
Krydskorrelationen giver en forskydning (lag) i samples mellem to
udsnit. Da udsnittene fra de to noder ikke starter på samme tidspunkt,
er tidsforskellen:

\[\Delta t = (t_{B} - t_{A}) + \frac{\mathit{lag}}{48000\ \mathrm{Hz}}\]

Her er \(t_{A}\) og \(t_{B}\) timestamps for de to udsnits første frame.

\paragraph{Små noter}

\begin{itemize}
\item
  Hændelsesdetekteringen kører i en separat tråd i hver af de to
  mikrofonnoder (de to slaver)
\item
  Mikrofonmodellen er ``SPH0645LM4H-B''
\item
  En mikrofonnode registrerer kun en hændelse, hvis begge dens
  mikrofoner har overskredet tærsklen
\item
  En hændelse behandles kun hvis begge slaver har detekteret en
  hændelse. Dette vurderes i master RPI.
\item
  Vi sampler med 48 kHz i ALSA-formatet S32\_LE
\item
  Hver mikrofonnodes to mikrofoner sidder i nabohjørner af målområdet
\item
  Timestamps leveres sammen med data fra ringbufferen og antages at være
  synkroniseret mellem noderne via PTP
\item
  Hver hændelse sendes som et udsnit på 50 ms (15 ms før og 35 ms efter
  triggerpunktet) til master over TCP
\end{itemize}

\end{document}
